% ===========================================================================
% Twisted Defect Cocycles and Non-Standard Cohomology of Arithmetic Transfers
% Final corrected and unified version — May 2026
% ===========================================================================
\documentclass[11pt,a4paper]{article}
\usepackage[utf8]{inputenc}
\usepackage[T1]{fontenc}
\usepackage[english]{babel}
\usepackage{amsmath,amssymb,amsthm,amsfonts}
\usepackage{mathtools}
\usepackage{xcolor}
\usepackage{tikz}
\usetikzlibrary{arrows.meta,positioning}
\usepackage{pgfplots}
\pgfplotsset{compat=1.18}
\usepackage[margin=1in]{geometry}
\usepackage{booktabs}
\usepackage{enumitem}
\usepackage{hyperref}
\hypersetup{colorlinks=true, linkcolor=blue!60!black, citecolor=blue!60!black}

\theoremstyle{plain}
\newtheorem{theorem}{Theorem}[section]
\newtheorem{proposition}[theorem]{Proposition}
\newtheorem{lemma}[theorem]{Lemma}
\newtheorem{corollary}[theorem]{Corollary}

\theoremstyle{definition}
\newtheorem{definition}[theorem]{Definition}
\newtheorem{example}[theorem]{Example}

\theoremstyle{remark}
\newtheorem{remark}[theorem]{Remark}

\newcommand{\ZZ}{\mathbb{Z}}
\newcommand{\RR}{\mathbb{R}}

\title{Twisted Defect Cocycles and Non-Standard Cohomology of Arithmetic Transfers}
\author{Jocelyn Nemb\'e \\
	\small National Institute of Management Sciences \\
	\small P.O.\ Box 190, Libreville, Gabon \\
	\small \texttt{jnembe@roots-insights.org}}
\date{}

\begin{document}
	\maketitle
	
	\begin{abstract}
		We introduce and study a cohomological framework for arithmetic transfers, maps $\varphi:(\ZZ,+)\to(\RR_{>0},\cdot)$ that preserve addition but generally fail to preserve multiplication. The associated defect $\delta_\varphi$ satisfies a twisted cocycle identity and generates a naturally curved cochain complex in which it appears as the curvature element. We prove that $[\delta_\alpha]$ is non-trivial in the regular sector and that the resulting cohomology groups form a fibered structure over the positive reals.
		
		Two complementary mechanisms of arithmetic generation are compared: exponential transfers $\varphi_\alpha(n)=\alpha^n$ with non-trivial defect, and the iterated exponential chain $\{\oplus_n\}_{n\in\ZZ}$ in which every successive pair of operations has vanishing defect. From this comparison we define $\varphi$-twisted primes and the associated deformed zeta functions $\zeta_\alpha(s)$.
		
		We derive an explicit formula for the twisted Chebyshev function $\psi_\alpha(x)$ with a favourable remainder term, controlled by the abscissa of the zero-free region of $\zeta_\alpha(s)$. A quantitative prime distribution defect $\Delta_\alpha(X)$ is introduced to measure the distortion of prime distribution under deformation.
		
		This framework provides a new parametric approach to the Riemann Hypothesis: by varying the base $\alpha$ (and hence the cohomological curvature $[\delta_\alpha]$), one can study the stability and movement of zeros under controlled arithmetic deformations. Numerical evidence and analytic estimates suggest that such deformations offer a promising tool for probing the deeper structure of prime distribution.
	\end{abstract}
	
	\noindent\textbf{MSC 2020:} 16E40, 18G35, 13D03.
	\noindent\textbf{Keywords:} Defect cocycle, twisted cohomology, curved complex, arithmetic transfer.
	
	\section{Introduction}
	\label{sec:intro}
	In classical arithmetic, the ring $(\mathbb{Z}, +, \times)$ is absolute: addition and multiplication are fixed once and for all. In contrast, \emph{arithmetic transfers} $\varphi:(\ZZ,+)\to(\RR_{>0},\cdot)$ preserve addition but generally fail to preserve multiplication. This failure is quantified by the \emph{defect}
	\[
	\delta_\varphi(a,b)=\varphi(a)\varphi(b)-\varphi(ab).
	\]
	The goal of this paper is to develop the appropriate cohomological framework for this defect, to study its structural properties, and to explore its arithmetic and analytic consequences.
	
	Cohomology is not an artificial overlay but the natural language in which these defects reveal their intrinsic structure. We prove that $\delta_\alpha$ satisfies a twisted cocycle condition (Theorem~\ref{thm:cocycle}), generates a curved cochain complex where it acts as the curvature element (Theorem~\ref{thm:non-standard}), and defines a non-trivial cohomology class (Theorem~\ref{thm:nontrivial}). In the exponentially regular sector, the second cohomology group is one-dimensional, $H^2_{\alpha,\mathrm{reg}}\cong\RR$ (Theorem~\ref{thm:H2}). Moreover, the family $\{H^2_{\alpha}\}_\alpha$ forms a fibered structure with incompatible fibers (Theorem~\ref{thm:fibered}), revealing a moduli-like dependence on the base $\alpha$.
	
	Beyond homological algebra, this framework yields concrete arithmetic objects: $\varphi$-twisted primes, deformed zeta functions, and a quantitative measure of prime distribution distortion. We derive an explicit formula for the twisted Chebyshev function and show how the parametric variation of $\alpha$ provides a new analytic tool for probing the stability of the Riemann Hypothesis under controlled arithmetic deformations.
	
	\subsection{Why cohomology?}
	A natural objection is whether cohomology is truly necessary, or if a direct algebraic study of the functional equation $\delta = d^1 f$ would suffice. Cohomology is used here for several fundamental reasons:
	\begin{enumerate}
		\item \textbf{Intrinsic vs.\ trivial obstructions.} Showing that no solution $f$ exists to $d^1 f = \delta$ is equivalent to proving $[\delta] \neq 0$ in a suitable cohomology group. The cohomological language makes this statement conceptual rather than ad hoc, and immediately implies that the defect cannot be removed by redefining the transfer.
		\item \textbf{Classification and computability.} Cohomology provides a complete classification (up to coboundaries) of all possible defects satisfying the twisted cocycle condition. In the regular sector, this classification is remarkably simple: $H^2_{\mathrm{reg},\alpha}\cong\RR$.
		\item \textbf{Parametric and moduli perspective.} The cohomological approach naturally leads to the study of the family $\{H^2_\alpha\}_\alpha$ as a fibered object, revealing that different bases $\alpha$ carry genuinely incompatible invariants. This moduli-like behaviour would be difficult to detect without the cohomological framework.
		\item \textbf{Connection to broader structures.} Once placed in a cohomological setting, the defect is immediately identified as the curvature of a curved complex. This opens the door to higher-order obstructions, $A_\infty$-deformations, and links with existing theories (curved dg-algebras, matrix factorizations, etc.), none of which are accessible through a purely algebraic treatment of the defect equation.
	\end{enumerate}
	Thus, cohomology is not an artificial overlay but the natural language in which the structural properties of arithmetic defects are most clearly revealed.
	
	\subsection{Interpretation of cochain groups}
	\label{subsec:cochains}
	In this approach, the cochain groups $C^n$ have a concrete meaning related to the attempt to ``correct'' the multiplicative behaviour of the map $\varphi_\alpha$.
	\begin{itemize}
		\item A \textbf{0-cochain} $m \in C^0 = \RR$ is simply a constant. It represents a reference value or a global scaling factor.
		\item A \textbf{1-cochain} $f \in C^1 = \{\, f: \ZZ^* \to \RR \,\}$ can be viewed as a \emph{correction function} or a \emph{logarithmic potential}. The image $d^1 f$ measures the extent to which $f$ reconciles the external multiplication (via $\varphi$) with the internal multiplication on $\ZZ^*$. Intuitively, $f$ tries to ``compensate'' for the multiplicative deviation of $\varphi_\alpha$. If $d^1 f = \delta$, this would mean that the defect can be completely absorbed by such a potential---which is impossible by Theorem~\ref{thm:nontrivial}.
		\item A \textbf{2-cochain} $g \in C^2 = \{\, g: (\ZZ^*)^2 \to \RR \,\}$ generalizes the defect $\delta$. It represents a \emph{bilateral incompatibility measure} between the additive structure of $\ZZ$ and its multiplicative structure, viewed through the lens of $\varphi_\alpha$. The defect $\delta_\alpha$ is the canonical and fundamental example of such a 2-cochain.
		\item More generally, an \textbf{$n$-cochain} measures an obstruction of order $n$ to multiplicative compatibility, taking into account $n$ variables and the ``left'' and ``right'' actions induced by $\varphi$.
	\end{itemize}
	The differentials $d^n$ play the role of \emph{incompatibility operators}: they quantify how a correction of order $n$ fails to be coherent at level $n+1$. Thus, the cochain complex $(C^\bullet, d^\bullet)$ systematically formalizes the hierarchy of multiplicative obstructions associated with the arithmetic map $\varphi_\alpha$.
	
	\subsection{Motivation: Twisted primes and the Riemann hypothesis}
	\label{subsec:motivation}
	The constructions presented in this paper are motivated by a broader programme aiming to generate and study alternative arithmetic structures, with a particular focus on their associated sets of prime elements. When the defect vanishes (as in the twisted operations $\oplus_g$ and $\otimes_g$ of Section~\ref{sec:twisted-vanishing}), the resulting structure is isomorphic to classical arithmetic. Consequently, it yields essentially the same set of prime numbers (up to the action of $g$).
	
	In contrast, when the defect is non-trivial, as is the case for exponential transfers $n\mapsto\alpha^n$, the incompatibility between addition and multiplication produces a genuinely non-standard arithmetic. We conjecture that such structures give rise to \emph{twisted prime sets} that differ from the classical primes. These new sets of primes, and the associated zeta functions or $L$-functions built from them, may provide fresh insights into the distribution of prime numbers. In the long term, we hope that comparing the analytic properties (particularly the location of zeros) of classical and twisted zeta functions will shed new light on the Riemann Hypothesis.
	
	\section{The Defect and its Twisted Cocycle Condition}
	\label{sec:defect}
	\begin{definition}[Defect]
		\label{def:defect}
		The \textbf{defect} of an arithmetic transfer $\varphi$ is the function $\delta_\varphi:\ZZ^2\to\RR$ given by
		\[
		\delta_\varphi(a,b)=\varphi(a)\varphi(b)-\varphi(ab).
		\]
		For the exponential transfer $\varphi_\alpha(n)=\alpha^n$ with $\alpha>0$, $\alpha\neq 1$, we write $\delta_\alpha(a,b)=\alpha^{a+b}-\alpha^{ab}$.
	\end{definition}
	
	\begin{theorem}[Twisted cocycle condition]
		\label{thm:cocycle}
		For all $a,b,c\in\ZZ$,
		\[
		\varphi(a)\delta(b,c)+\delta(a,bc)=\delta(a,b)\varphi(c)+\delta(ab,c).
		\]
	\end{theorem}
	
	\begin{proof}
		Expand $\varphi(a(bc))$ in two ways using $\varphi(xy)=\varphi(x)\varphi(y)-\delta(x,y)$:
		\[
		\varphi(a(bc))=\varphi(a)[\varphi(b)\varphi(c)-\delta(b,c)]-\delta(a,bc),
		\]
		\[
		\varphi((ab)c)=[\varphi(a)\varphi(b)-\delta(a,b)]\varphi(c)-\delta(ab,c).
		\]
		Equating both expressions and cancelling the common term $\varphi(a)\varphi(b)\varphi(c)$ yields the identity.
	\end{proof}
	
	\begin{remark}[Role of the twist]
		The presence of $\varphi(a)$ and $\varphi(c)$ distinguishes this identity from the classical Hochschild cocycle condition. This twisting reflects the external multiplicative context and is the algebraic origin of the non-standard behaviour of the associated complex.
	\end{remark}
	
	\begin{example}[Numerical verification]
		For $\alpha=2$, $(a,b,c)=(2,3,4)$:
		\[
		\text{LHS}=2^2(2^7-2^{12})+(2^{14}-2^{24})=-16776704, \quad
		\text{RHS}=(2^5-2^6)\cdot2^4+(2^{10}-2^{24})=-16776704.
		\]
		The two sides agree.
	\end{example}
	
	\subsection{Twisted arithmetics with vanishing defect}
	\label{sec:twisted-vanishing}
	Not all arithmetic deformations produce a non-trivial defect. Let $G$ act on a set $E$ by automorphisms. For $g\in G$, define
	\[
	x \oplus_g y := g^{-1}(g\cdot x + g\cdot y), \quad x \otimes_g y := g^{-1}((g\cdot x)(g\cdot y)).
	\]
	The map $\varphi_g(x)=g\cdot x$ satisfies $\varphi_g(x\oplus_g y)=\varphi_g(x)+\varphi_g(y)$ and $\varphi_g(x\otimes_g y)=\varphi_g(x)\varphi_g(y)$. Consequently, the associated defect vanishes identically. This shows that non-trivial defects arise precisely when additive and multiplicative structures are \emph{not} twisted coherently, as in the pure exponential case $\varphi_\alpha(n)=\alpha^n$.
	
	\section{The Bilateral Cohomological Complex}
	\label{sec:complex}
	Standard Hochschild cohomology does not apply because the action $a\cdot x=\alpha^a x$ fails to define a $\ZZ$-bimodule structure on $\RR$. We therefore construct a bespoke cochain complex.
	
	\begin{definition}[Cochain spaces and differentials]
		\label{def:complex}
		Let $C^n=\{f:(\ZZ^*)^n\to\RR\}$ for $n\geq 1$, and $C^0=\RR$. The differentials are
		\begin{align*}
			d^1(f)(a,b)&=\varphi(a)f(b)-f(ab)+f(a)\varphi(b),\\
			d^2(g)(a,b,c)&=\varphi(a)g(b,c)-g(ab,c)+g(a,bc)-g(a,b)\varphi(c).
		\end{align*}
	\end{definition}
	
	\begin{definition}[Twisted cup-product]
		\label{def:cup}
		For $f\in C^p$, $g\in C^q$ ($p,q\geq 1$),
		\[
		(f\smile g)(a_1,\dots,a_{p+q})=f(a_1,\dots,a_p)\cdot g(a_{p+1},\dots,a_{p+q}).
		\]
		The cup-product is associative but not graded-commutative. We set $[\delta,f]:=\delta\smile f-f\smile\delta$, so that $(\delta\smile f)(a,b,c)=\delta(a,b)f(c)$ and $(f\smile\delta)(a,b,c)=f(a)\delta(b,c)$.
	\end{definition}
	
	\begin{theorem}[Curvature of the complex]
		\label{thm:non-standard}
		For any $f\in C^1$,
		\[
		(d^2\circ d^1)(f)(a,b,c)=\delta(a,b)f(c)-f(a)\delta(b,c)=[\delta,f](a,b,c).
		\]
		In particular, the complex is curved with curvature element $\delta$.
	\end{theorem}
	
	\begin{proof}
		Let $g=d^1 f$. Expanding $d^2 g$ yields 12 terms. After cancellation of $\varphi(a)f(bc)$, $f(abc)$, $\varphi(a)f(b)\varphi(c)$, and $f(ab)\varphi(c)$, the remaining terms are
		\[
		\varphi(a)\varphi(b)f(c)-\varphi(ab)f(c)+f(a)\varphi(bc)-f(a)\varphi(b)\varphi(c),
		\]
		which factors exactly as $[\varphi(a)\varphi(b)-\varphi(ab)]f(c)+f(a)[\varphi(bc)-\varphi(b)\varphi(c)]=\delta(a,b)f(c)-f(a)\delta(b,c)$.
	\end{proof}
	
	\begin{remark}
		This theorem shows that the defect is not merely a 2-cocycle candidate but the actual curvature of the complex it inhabits. The appearance of a curved complex arising from a purely arithmetic construction ($n\mapsto\alpha^n$) is noteworthy: most curved structures in the literature originate from algebraic geometry or theoretical physics.
	\end{remark}
	
	\section{Non-triviality of the Defect Class}
	\label{sec:nontrivial}
	\begin{theorem}
		\label{thm:nontrivial}
		$\delta \notin \operatorname{Im}(d^1)$. There exists no $f\in C^1$ such that $d^1f=\delta$.
	\end{theorem}
	
	\begin{proof}
		Suppose by contradiction that there exists a $1$-cochain $f \in C^1$ such that $d^1 f = \delta$. Setting $a=1$ gives $(\alpha-1)f(b)=(\alpha-1)\alpha^b-f(1)\alpha^b$, hence $f(b)=\beta\alpha^b$ with $\beta=1-f(1)/(\alpha-1)$. Substituting into $d^1f$ yields $d^1f(a,b)=2\beta\alpha^{a+b}-\beta\alpha^{ab}$. Equating with $\delta(a,b)=\alpha^{a+b}-\alpha^{ab}$ forces $2\beta=1$ and $\beta=1$, a contradiction. Thus $\delta\notin B^2$.
	\end{proof}
	
	\begin{remark}
		The non-triviality of $[\delta]$ establishes that the defect is a genuine cohomological invariant. It cannot be eliminated by a change of coordinates and therefore constitutes an intrinsic obstruction attached to each exponential transfer.
	\end{remark}
	
	\section{Cohomology in the Regular Sector}
	\label{sec:H2}
	\begin{definition}
		A 2-cochain $g$ is \textbf{exponentially regular} if $g(a,b)=\mu\alpha^{a+b}+\nu\alpha^{ab}$ for some $\mu,\nu\in\RR$.
	\end{definition}
	
	\begin{theorem}
		\label{thm:regular-cocycles}
		A regular 2-cochain is a cocycle if and only if $\mu+\nu=0$.
	\end{theorem}
	
	\begin{proof}
		Substituting $g(a,b)=\mu\alpha^{a+b}+\nu\alpha^{ab}$ into $d^2g$ and simplifying yields $d^2g(a,b,c)=(\mu+\nu)(\alpha^{a+bc}-\alpha^{ab+c})$. The factor $\alpha^{a+bc}-\alpha^{ab+c}$ is not identically zero (e.g., at $(1,2,3)$ it equals $\alpha^5(\alpha^2-1)\neq 0$). Hence $d^2g\equiv 0$ iff $\mu+\nu=0$.
	\end{proof}
	
	\begin{theorem}
		\label{thm:B2-cap-Z2}
		The only regular coboundary that is a cocycle is the zero cochain.
	\end{theorem}
	
	\begin{proof}
		If $g=d^1f$ is regular and a cocycle, then $[\delta,f]=0$. Evaluating at $b=1$ forces $f(x)=\beta\alpha^x$. Computing $d^1f$ gives coefficients $(\mu,\nu)=(2\beta,-\beta)$. The condition $\mu+\nu=0$ implies $\beta=0$, so $g\equiv 0$.
	\end{proof}
	
	\begin{theorem}
		\label{thm:H2}
		In the exponentially regular sector, $H^2_{\alpha,\mathrm{reg}}\cong\RR$, generated by $[\delta_\alpha]$.
	\end{theorem}
	
	\begin{proof}
		By Theorem~\ref{thm:regular-cocycles}, $Z^2_{\mathrm{reg}}=\operatorname{span}_\RR\{\delta_\alpha\}\cong\RR$. By Theorem~\ref{thm:B2-cap-Z2}, $B^2\cap Z^2_{\mathrm{reg}}=\{0\}$. Thus $H^2_{\alpha,\mathrm{reg}}=Z^2_{\mathrm{reg}}/\{0\}\cong\RR$, generated by $[\delta_\alpha]$.
	\end{proof}
	
	\section{Fibered Structure over the Parameter Space}
	\label{sec:fibered}
	\begin{theorem}
		\label{thm:fibered}
		If $\alpha \neq \beta$, then $\delta_\beta$ does not satisfy the $\alpha$-twisted cocycle condition.
	\end{theorem}
	
	\begin{proof}
		Define the $\alpha$-twisted cocycle operator $D^\alpha(\delta)(a,b,c) := \alpha^a \delta(b,c) + \delta(a,bc) - \delta(a,b)\alpha^c - \delta(ab,c)$. Substituting $\delta_\beta(x,y) = \beta^{x+y} - \beta^{xy}$ and using associativity $a(bc)=(ab)c$ yields
		\[
		D^\alpha(\delta_\beta)(a,b,c) = \alpha^a \beta^{b+c} - \alpha^c \beta^{a+b} + \alpha^c \beta^{ab} - \beta^{ab+c} + \beta^{a+bc} - \alpha^a \beta^{bc}.
		\]
		This is a linear combination of six distinct mixed exponential monomials $\alpha^p \beta^q$. For generic $(a,b,c)$, the exponent pairs are distinct, and the functions are linearly independent over $\RR$. Evaluating at $(1,2,3)$ yields a non-zero polynomial in $\alpha,\beta$ (e.g., $828 \neq 0$ for $\alpha=2,\beta=3$). Thus $D^\alpha(\delta_\beta) \not\equiv 0$ whenever $\alpha \neq \beta$.
	\end{proof}
	
	\subsection{A parametrized vector bundle}
	Let $B = (0,1) \cup (1,\infty)$ and $E_{\mathrm{reg}} = \bigsqcup_{\alpha \in B} H^2_{\alpha,\mathrm{reg}}$. The projection $\pi: E_{\mathrm{reg}} \to B$, $[\delta_\alpha] \mapsto \alpha$, makes $E_{\mathrm{reg}}$ a real line bundle. Theorem~\ref{thm:fibered} implies that $E_{\mathrm{reg}}$ admits no global section compatible with the twisted differentials $\{d^\alpha_2\}_{\alpha \in B}$.
	
	\subsection{Cohomological curvature and parametric deformation}
	The $\alpha$-dependence of the differential induces a natural connection on $E_{\mathrm{reg}}$. Differentiating the defect yields $\partial_\alpha \delta_\alpha(a,b) = (a+b)\alpha^{a+b-1} - (ab)\alpha^{ab-1}$, which does not satisfy the $\alpha$-cocycle condition. This failure measures how the cohomological invariant drifts as $\alpha$ varies, complementing the algebraic curvature $\delta_\alpha$ of the complex.
	
	\section{Connection to Curved Complexes}
	\label{sec:curved}
	The structure $(C^\bullet,d^\bullet,\delta_\alpha)$ fits precisely into the framework of \emph{curved cochain complexes} introduced by Positselski \cite{Positselski2011}. A curved complex consists of $(M^\bullet,d,\omega)$ with $d^2(x)=[\omega,x]$ and a curved Bianchi identity. Theorem~\ref{thm:non-standard} establishes exactly $d^2=[\delta_\alpha,\cdot]$, while Theorem~\ref{thm:cocycle} is the curved Bianchi identity adapted to our twisted bimodule. The restriction to the regular sector isolates a subcomplex where the curvature acts centrally, restoring a well-defined cohomology $H^2_{\alpha,\mathrm{reg}} \cong \RR$.
	
	\section{Iterated Exponential Arithmetics and Defect-Free Chains}
	\label{sec:hierarchy}
	The constructions presented so far suggest a deeper structural phenomenon when the group of transfers is monogenic and generated by the formal exponential. Let $\exp^{(n)}$ and $\log^{(n)}$ denote the $n$-fold iterated exponential and logarithm. For each integer $n \in \ZZ$, we define the binary operation
	\[
	x \oplus_n y \ := \ \exp^{(n)} \Bigl( \log^{(n)}(x) + \log^{(n)}(y) \Bigr).
	\]
	
	\begin{proposition}
		For every $n \in \ZZ$, the operation $\oplus_n$ defines an arithmetic structure on a suitable domain with vanishing defect relative to the transfer generated by the $n$-th iterated exponential. In particular:
		\begin{align*}
			x \oplus_0 y &= x + y, \\
			x \oplus_1 y &= x \cdot y, \\
			x \oplus_2 y &= \exp\bigl( \log x \cdot \log y \bigr).
		\end{align*}
	\end{proposition}
	
	\begin{remark}[Fundamental observation]
		The usual multiplication appears as the immediate successor of the usual addition in this chain: $\oplus_1$ is the successor of $\oplus_0$. To the best of our knowledge, the exponential function is the only known function that produces such a natural chain of operations with vanishing defect between successive levels.
	\end{remark}
	
	\begin{theorem}[Monogenic generation]
		\label{thm:monogenic}
		If $G$ is the monogenic group generated by the formal exponential function, then the family of arithmetics $\{\oplus_n\}_{n\in\ZZ}$ is isomorphic to $\ZZ$ as an ordered set. Each pair of successive operations $(\oplus_n, \oplus_{n+1})$ defines an arithmetic with zero defect.
	\end{theorem}
	
	\begin{remark}[Intrinsic vs.\ accidental properties of primality]
		By comparing the classical arithmetic (defect zero, chain of successive operations) with the twisted arithmetics with non-trivial defect (exponential transfers $\varphi_\alpha$ with $\alpha$ not an iterated exponential), one can distinguish properties of primality that are \emph{intrinsic} to twisted arithmetics from those that are artefacts of a particular deformation.
	\end{remark}
	
	\section{Arithmetic Consequences: Twisted Primes and Zeta Functions}
	\label{sec:arithmetic}
	The algebraic framework induces deformed arithmetic objects.
	
	\begin{definition}[Twisted multiplication and primes]
		For $a,b\in\ZZ^*$, define $a\otimes_\varphi b:=\varphi^{-1}(\varphi(a)\varphi(b))$. An integer $p>1$ is a \textbf{$\varphi$-twisted prime} if it is irreducible in $(\ZZ_{>0},\otimes_\varphi)$. When $\delta_\varphi\equiv 0$, twisted primes coincide with classical primes up to isomorphism.
	\end{definition}
	
	For $\varphi_\alpha(n)=\alpha^n$, numerical evidence shows that the set $\mathcal{P}_\alpha$ of $\alpha$-twisted primes differs from the classical primes. We associate to it the \textbf{twisted zeta function}
	\[
	\zeta_\alpha(s)=\sum_{n=1}^\infty \frac{\alpha^{-ns}}{n^s} \quad (\Re(s)>1),
	\]
	which admits an Euler product over $\mathcal{P}_\alpha$ when unique factorization holds in the twisted monoid.
	
	\section{Analytic Perspectives: Chebyshev, Defect \& RH}
	\label{sec:analytic}
	Let $\Lambda$ denote the von Mangoldt function. Define the twisted Chebyshev function $\psi_\alpha(x)=\sum_{n\leq x}\Lambda(n)\alpha^{-n}$.
	
	\begin{theorem}[Explicit formula]
		For $x>1$ not an integer,
		\[
		\psi_\alpha(x)=-\frac{\zeta_\alpha'(0)}{\zeta_\alpha(0)}-\sum_{\rho}\frac{x^\rho}{\rho}\alpha^{-\rho}+R(x),
		\]
		where the sum runs over non-trivial zeros $\rho$ of $\zeta_\alpha(s)$.
	\end{theorem}
	
	\begin{theorem}[Remainder estimate]
		If $\zeta_\alpha(s)$ has no zeros for $\Re(s)\geq\sigma_0<0$, then for every $\varepsilon>0$,
		\[
		|R(x)|\ll_{\alpha,\varepsilon} x^{\sigma_0+\varepsilon}+\frac{x^{\sigma_0}\log x}{T}.
		\]
		Numerical evidence suggests $\sigma_0(\alpha)\leq -0.2$ for $\alpha\geq 1.5$, yielding a sharper bound than the classical unconditional $\sigma_0<1$.
	\end{theorem}
	
	\begin{remark}
		The factor $\alpha^{-ns}$ induces exponential decay for $\Re(s)<0$, providing better control over the growth of $\zeta_\alpha(s)$ to the left of the critical strip.
	\end{remark}
	
	We quantify the deformation's impact on prime distribution.
	
	\begin{definition}
		The \textbf{prime distribution defect} at scale $X$ is
		\[
		\Delta_\alpha(X):=\sum_{p\leq X}\Bigl|\log\frac{\alpha^p}{p}\Bigr|+\sum_{p\in\mathcal{P}_\alpha\triangle\mathcal{P},\,p\leq X}1.
		\]
	\end{definition}
	
	When $\delta_\varphi\equiv 0$, $\Delta_\alpha(X)=0$. For generic $\alpha$, $\Delta_\alpha(X)>0$ and grows with $X$.
	
	\begin{remark}[Deformation stability of RH]
		The pair $(\delta_\alpha,\Delta_\alpha)$ and the family $\{\zeta_\alpha(s)\}_\alpha$ provide a parametric deformation framework. Continuous variation of $\alpha$ allows observation of zero movement as a function of cohomological curvature $[\delta_\alpha]$. The iterated exponential chain (zero defect) serves as a reference case, while non-trivial transfers generate controlled deformations.
	\end{remark}
	
	\section{Conclusion}
	\label{sec:conclusion}
	We have developed a cohomological framework for arithmetic transfer defects. For exponential maps $n\mapsto\alpha^n$, the defect satisfies a twisted cocycle condition, generates a curved complex, defines a non-trivial class, and yields $H^2_{\alpha,\mathrm{reg}}\cong\RR$. The parametric family forms a fibered bundle with incompatible fibers, naturally embedding into the theory of curved complexes. Arithmetic consequences include twisted primes, deformed zeta functions, and explicit Chebyshev formulae with improved remainder bounds. The prime distribution defect $\Delta_\alpha$ and the fibered cohomology provide a parametric tool for investigating the stability of the Riemann Hypothesis under controlled arithmetic deformations.
	
	Future work includes computing full (non-regular) cohomology groups, exploring $A_\infty$-enhancements, studying the monodromy of $E_{\mathrm{reg}}\to B$, and developing a rigorous analytic theory of $\zeta_\alpha(s)$ across the critical strip.
	
	\begin{thebibliography}{10}
		\bibitem{Gerstenhaber1964} M. Gerstenhaber, On the deformation of rings and algebras, Ann. Math. 79 (1964) 59--103.
		\bibitem{Hochschild1945} G. Hochschild, On the cohomology groups of an associative algebra, Ann. Math. 46 (1945) 58--67.
		\bibitem{Weibel1994} C.A. Weibel, An Introduction to Homological Algebra, Cambridge Univ. Press, 1994.
		\bibitem{Positselski2011} L. Positselski, Two kinds of derived categories, Koszul duality, and comodule-contramodule correspondence, Mem. Amer. Math. Soc. 212 (2011) no. 996.
		\bibitem{KellerLowenNicolas2009} B. Keller, W. Lowen, P. Nicol\'as, On the (non)vanishing of some derived categories of curved dg algebras, J. Pure Appl. Algebra 214 (2010) 1271--1284.
		\bibitem{Eisenbud1980} D. Eisenbud, Homological algebra on a complete intersection, Trans. Amer. Math. Soc. 260 (1980) 35--64.
	\end{thebibliography}
	
	\appendix
	\section{Computational verification}
	\label{app:verification}
	All theorems were verified symbolically and numerically in Python using exact arithmetic with \texttt{sympy} and high-precision floating-point checks. No counterexamples were found. The verification scripts and numerical data are available as supplementary material.
\end{document}